\documentclass[11pt]{article}

\usepackage[T1]{fontenc}
\usepackage[utf8]{inputenc}
\usepackage{lmodern}
\usepackage[margin=1in]{geometry}
\usepackage{amsmath,amssymb,amsthm,mathtools}
\usepackage{booktabs,array}
\usepackage{xcolor}
\usepackage{hyperref}
\usepackage{xurl}
\usepackage{microtype}
\newcommand{\OPHPaperReleaseID}{r2011}
\newcommand{\OPHPaperReleaseDate}{August 1, 2026}
\newcommand{\OPHPaperReleaseBanner}{%
\begin{center}
\small\textbf{Paper release:} \texttt{\OPHPaperReleaseID}\quad\textbf{Released:} \OPHPaperReleaseDate
\end{center}
}
\newcommand{\PaperReleaseID}{\OPHPaperReleaseID}
\newcommand{\PaperReleaseDate}{\OPHPaperReleaseDate}
\newcommand{\PaperReleaseBanner}{%
\begin{center}
\small\textbf{Paper release:} \texttt{\PaperReleaseID}\quad\textbf{Released:} \PaperReleaseDate
\end{center}
}


\hypersetup{
  colorlinks=true,
  linkcolor=blue,
  citecolor=blue,
  urlcolor=blue,
  pdftitle={The Positive-Chamber Koide Identity for Icosahedral Face Circulants},
  pdfauthor={Bernhard Mueller}
}
\setlength{\parindent}{0pt}
\setlength{\parskip}{0.55em}
\setlength{\emergencystretch}{3em}
\renewcommand{\arraystretch}{1.12}

\newtheorem{theorem}{Theorem}
\newtheorem{proposition}[theorem]{Proposition}
\theoremstyle{remark}
\newtheorem{remark}[theorem]{Remark}

\newcommand{\C}{\mathbb C}
\newcommand{\R}{\mathbb R}
\newcommand{\one}{\mathbf 1}
\newcommand{\Tr}{\operatorname{Tr}}

\title{\textbf{The Positive-Chamber Koide Identity\\
for Icosahedral Face Circulants}}
\author{Bernhard Mueller\thanks{Corresponding author:
\href{mailto:bernhard@floatingpragma.ai}{\texttt{bernhard@floatingpragma.ai}}}\\
Pragma Research Inc., Washington, United States}
\date{\OPHPaperReleaseDate}

\begin{document}
\maketitle
\OPHPaperReleaseBanner

\begin{abstract}
The stabilizer of an oriented icosahedral face acts regularly on its three
corners. Every Hermitian response commuting with this cyclic action is a
circulant
\(C=aI+bR+\overline bR^2\), \(R^3=I\). In the chamber where all three
eigenvalues are nonnegative and may be read as square-root masses, its Koide
invariant is
\[
Q=\frac13+\frac23\left(\frac{|b|}{a}\right)^2.
\]
Thus \(Q=2/3\) is exactly equivalent to the single balance condition
\(|b|/a=1/\sqrt2\). The phase of \(b\) drops out of \(Q\) and continues to
control the two mass ratios jointly. A finite tracial
Gelfand--Naimark--Segal construction with equal rank-two event blocks gives
the balanced modulus as a conditional theorem. The local face fiber has no
constructed physical chiral-family attachment, so the result is a
representation-theoretic identity and a finite conditional implication,
without a charged-lepton mass prediction. The numerical response model
associated with this construction uses target-informed inputs. Its
near-equality with the Particle Data Group central Koide coordinate is a
comparison-only diagnostic.
\end{abstract}

\noindent\textbf{Keywords:} Koide relation, circulant matrices, cyclic
symmetry, icosahedral symmetry, finite event algebra, GNS representation

\section{The face-corner carrier}

Koide's charged-lepton relation~\cite{koide1983} is usually written in terms
of the three positive square-root masses. The construction below identifies
the exact condition it imposes on a cyclic Hermitian response.

Let \(G=A_5\), the alternating group on five letters, act by proper rotations
on the icosahedron. The twenty outward
oriented faces form the transitive orbit \(G/C_3\). The stabilizer of one face
cyclically permutes its three corners, so every face carries a local copy of
the regular \(C_3\) representation. If \(R\) denotes the cyclic shift, then
\[
R^3=I,\qquad R^\dagger=R^2.
\]
The commutant of this regular action is the three-dimensional circulant
algebra
\[
\{x_0I+x_1R+x_2R^2:x_j\in\C\}.
\]
Hermiticity restricts a response to
\begin{equation}
C=aI+bR+\overline bR^2,\qquad a\in\R,\quad b\in\C.
\label{eq:circulant}
\end{equation}
Write \(b=\rho e^{i\delta}\), with \(\rho=|b|\geq0\). Fourier
diagonalization gives
\begin{equation}
\lambda_k=a+2\rho\cos\left(\delta+\frac{2\pi k}{3}\right),
\qquad k=0,1,2.
\label{eq:eigenvalues}
\end{equation}
The unordered spectrum is independent of the chosen face representative.
This is a local bundle statement. The sixty face-corner flags form the
regular \(A_5\) torsor; the geometry alone supplies no canonical global
three-dimensional physical family space.

\section{The positive-chamber identity}

\begin{theorem}[Positive-chamber Koide identity]
\label{thm:koide}
Let \(C\) be the Hermitian circulant in
Eq.~\eqref{eq:circulant}, with \(a>0\). Suppose its three eigenvalues in
Eq.~\eqref{eq:eigenvalues} are nonnegative and, for one \(s>0\),
\[
\sqrt{m_k}=\sqrt{s}\lambda_k.
\]
Then
\begin{equation}
Q:=\frac{\sum_{k=0}^2m_k}
{\left(\sum_{k=0}^2\sqrt{m_k}\right)^2}
=\frac13+\frac23\left(\frac{\rho}{a}\right)^2.
\label{eq:koideformula}
\end{equation}
Consequently,
\begin{equation}
Q=\frac23
\quad\Longleftrightarrow\quad
\frac{\rho}{a}=\frac1{\sqrt2}.
\label{eq:balance}
\end{equation}
\end{theorem}

\begin{proof}
Set \(c_k=\cos(\delta+2\pi k/3)\). The roots-of-unity identities give
\[
\sum_{k=0}^2c_k=0,
\qquad
\sum_{k=0}^2c_k^2=\frac32.
\]
It follows that
\[
\sum_k\lambda_k=3a,
\qquad
\sum_k\lambda_k^2=3a^2+6\rho^2.
\]
The scale \(s\) cancels from \(Q\), yielding
Eq.~\eqref{eq:koideformula}. Since \(\rho/a\geq0\),
Eq.~\eqref{eq:koideformula} equals \(2/3\) exactly at
\(\rho/a=1/\sqrt2\).
\end{proof}

\begin{remark}[The chamber is part of the theorem]
The signed trace calculation holds algebraically for every \(\delta\).
Physical square roots use \(|\lambda_k|\) when an eigenvalue is negative, and
then the denominator is not the signed trace. At
\(\rho/a=1/\sqrt2\), the positive chamber is
\[
|\delta|\leq\frac{\pi}{12}\pmod{\frac{2\pi}{3}}.
\]
Theorem~\ref{thm:koide} makes no physical Koide statement outside that
chamber.
\end{remark}

\begin{remark}[What the phase contains]
Inside the positive chamber, \(Q\) is independent of \(\delta\). The three
eigenvalues, and hence the two reported mass ratios, vary with \(\delta\).
The balance condition removes one scale-free degree of freedom from a
three-mass spectrum. It does not determine the ratios.
\end{remark}

\section{Conditional finite tracial balance}

The balance condition in Eq.~\eqref{eq:balance} can arise exactly from a
finite event packet. This subsection states the packet as an independent
conditional theorem.

Let
\[
\mathcal V=\one\oplus\chi\oplus\overline\chi,
\qquad
\mathcal H_{\mathrm{or}}=\C^2,
\qquad
\mathcal A=B(\mathcal V\otimes\mathcal H_{\mathrm{or}})
\simeq M_6(\C).
\]
Here \(\one\) is the neutral cyclic mode, \(\chi\oplus\overline\chi\) is the
two-dimensional charged plane, and \(\mathcal H_{\mathrm{or}}\) is a
two-state orientation record. Let \(P_0\) and \(P_c\) project onto the neutral
line and charged plane. For a rank-one oriented event \(q_+\), define
\[
Z_0=P_0\otimes I_{\mathrm{or}},
\qquad
Z_c=P_c\otimes q_+,
\qquad
E_+=Z_0+Z_c.
\]
Both \(Z_0\) and \(Z_c\) have rank two.

\begin{theorem}[Conditional finite tracial-GNS balance]
\label{thm:gns}
Condition the normalized trace state \(I_6/6\) on \(E_+\). Map the orthonormal
cyclic basis of the resulting square-root amplitude to
\((I,R,R^2)\) in \(L^2(M_3(\C),\tau_3)\). Then
\[
p_0=p_c=\frac12,
\qquad
a=\sqrt{p_0},
\qquad
\rho=\sqrt{\frac{p_c}{2}},
\]
and therefore
\[
\frac{\rho}{a}=\frac1{\sqrt2},
\qquad
Q=\frac23
\]
whenever the positive square-root-mass reading of
Theorem~\ref{thm:koide} is supplied.
\end{theorem}

\begin{proof}
Normalized-trace conditioning assigns probability proportional to block
rank. Both accepted blocks have rank two, so \(p_0=p_c=1/2\). In
\(L^2(M_3(\C),\tau_3)\), the operators \(I,R,R^2\) are orthonormal. The
canonical square-root vector therefore has neutral coefficient
\(\sqrt{p_0}\) and two charged coefficients of modulus
\(\sqrt{p_c/2}\). Their ratio is \(1/\sqrt2\). Theorem~\ref{thm:koide}
then gives \(Q=2/3\) on the positive chamber.
\end{proof}

Theorem~\ref{thm:gns} is finite and exact. Its event algebra, block choice,
conditioned trace, and cyclic response map are premises. A physical
charged-lepton statement additionally requires a quotient-visible map from
the local face bundle to one chiral three-family response, preservation of
the relevant block powers, and a mass readout. Those objects are absent from
the theorem.

\section{Numerical diagnostic and provenance}

Using the Particle Data Group 2026 central masses
\cite{pdg2026},
\[
(m_e,m_\mu,m_\tau)
=(0.51099895069,\ 105.6583755,\ 1776.93)\ {\rm MeV},
\]
gives
\[
Q_{\mathrm{PDG}}=0.6666644634026367.
\]
The declared minimal complete public response model gives
\[
\begin{aligned}
Q_{\mathrm{MCPR}}&=0.6666644634090389,\\
\rho/a&=0.7071044442750720,\\
\frac{\rho/a-1/\sqrt2}{1/\sqrt2}
&=-3.3049\times10^{-6}.
\end{aligned}
\]
The response dimensions, path table, amplitude, phase, and determinant
exponent are target-informed inputs. The numerical response is therefore
neither blind nor source-derived. The proximity of the two displayed \(Q\)
values carries no significance or predictive weight.

\begin{center}
\begin{tabular}{@{}p{0.30\textwidth}p{0.58\textwidth}@{}}
\toprule
Statement & Scope \\
\midrule
Eq.~\eqref{eq:koideformula} & exact algebra in the nonnegative-eigenvalue chamber \\
Theorem~\ref{thm:gns} & exact implication conditional on the declared finite event packet \\
Icosahedral face fiber & exact local \(A_5/C_3\) bundle and unordered spectrum \\
Physical family attachment & work in progress \\
Phase and numerical ratios & work in progress \\
MCPR numerical proximity & target-informed comparison-only diagnostic \\
\bottomrule
\end{tabular}
\end{center}

\section{Why residual symmetry does not finish the ratios}

The remaining family-shape problem can be stated in the real
five-dimensional representation
\[
W_5\simeq\operatorname{Sym}^2_0(\R^3),
\]
the traceless symmetric \(3\times3\) matrices with action
\(A\mapsto gAg^T\).

\begin{proposition}[Residual-stabilizer boundary]
Invariance under a threefold or fivefold rotation about an axis \(n\) forces
\[
A=\alpha(nn^T-I/3),
\]
which has a double eigenvalue. Invariance under a twofold rotation has a
three-dimensional linear fixed locus, hence two projective parameters, and
that locus admits simple spectrum.
\end{proposition}

\begin{proof}
Choose the rotation axis as the third coordinate. Decompose a traceless
symmetric matrix into a planar scalar, planar spin-two anisotropy, transverse
vector, and axial entry. Rotation by \(\theta\) acts on the anisotropy by
\(2\theta\) and on the transverse vector by \(\theta\). For
\(\theta=2\pi/3\) or \(2\pi/5\), neither block has a nonzero fixed vector, so
only the axial combination remains. For \(\theta=\pi\), the planar symmetric
block survives:
\[
A=\begin{pmatrix}
u&v&0\\
v&w&0\\
0&0&-u-w
\end{pmatrix}.
\]
This space has dimension three and contains matrices with three distinct
eigenvalues.
\end{proof}

Threefold and fivefold symmetry cannot produce three distinct masses.
Twofold symmetry leaves two scale-free coordinates, exactly enough to carry
the two mass ratios without selecting them. A numerical spectrum therefore
requires a specific invariant potential derived from the screen dynamics.
This proposition locates the missing input; it supplies no replacement fit.

\section{Formal proof boundary}

The algebraic core has a Lean~4 formalization. It proves the root sum, square
sum, quotient formula, reciprocal-square-root identity, and the
nonnegative-modulus equivalence
\[
\frac13+\frac23r^2=\frac23
\quad\Longleftrightarrow\quad
r=\frac1{\sqrt2}.
\]
The formal module deliberately does not identify signed eigenvalues with
physical square roots. Positivity of all three eigenvalues, the phase
selection, and the source-to-charged-family attachment remain outside its
statement. This division matches Theorem~\ref{thm:koide}: the checked algebra
is exact, and the physical reading is a separate premise.

\section{Conclusion}

The regular three-corner face carrier turns the Koide relation into one
transparent modulus condition. In the positive chamber,
\[
Q=\frac23
\quad\Longleftrightarrow\quad
\rho/a=\frac1{\sqrt2}.
\]
The finite tracial event packet supplies this balance conditionally through
equal block weights and the canonical square-root representation. The result
does not determine the phase or mass ratios and does not attach the local
face fiber to physical charged leptons. Its durable content is an exact
circulant identity, an exact finite conditional balance theorem, and a sharp
statement of the remaining physical work.

\begin{thebibliography}{9}

\bibitem{koide1983}
Y.~Koide,
\emph{New View of Quark and Lepton Mass Hierarchy},
Physical Review D \textbf{28}, 252 (1983).
\url{https://doi.org/10.1103/PhysRevD.28.252}

\bibitem{pdg2026}
Particle Data Group,
\emph{Review of Particle Physics: 2026 particle listings},
2026.
\url{https://pdg.lbl.gov/2026/listings/particle_properties.html}

\end{thebibliography}

\end{document}
